\chapter*{CONCLUSION GÉNÉRALE}
\markboth{\MakeUppercase{CONCLUSION GÉNÉRALE}}{}
\addcontentsline{toc}{chapter}{CONCLUSION GÉNÉRALE}
\adjustmtc
\thispagestyle{MyStyle}

La réalisation de ce projet de conversion d'images braille en texte lisible a été une expérience enrichissante tant sur le plan technique que personnel. À travers le développement d'un modèle d'apprentissage automatique, nous avons exploré diverses approches et outils, tels que TensorFlow et Keras, pour traiter des données visuelles complexes et en extraire des informations pertinentes. 

Au cours de ce travail, nous avons pu comprendre les défis associés à la reconnaissance de caractères en braille et les solutions techniques permettant de surmonter ces obstacles. Le modèle entraîné a démontré une capacité prometteuse à convertir les images de braille en texte lisible, ce qui constitue un pas important vers l'accessibilité pour les personnes malvoyantes.

De plus, ce projet s'inscrit dans un cadre plus large d'innovation en matière de technologie d'impression braille. La collaboration avec l'équipe de TunBra-Msign a été essentielle pour orienter le développement du modèle et s'assurer qu'il réponde aux besoins spécifiques des utilisateurs finaux.

En conclusion, ce projet a non seulement renforcé mes compétences en apprentissage automatique et en traitement d'images, mais il a également renforcé mon engagement envers des solutions technologiques qui améliorent la vie des personnes en situation de handicap. Je suis convaincu que cette recherche pourra être prolongée et approfondie, ouvrant la voie à des applications encore plus avancées dans le domaine de l'accessibilité. À l'avenir, il serait intéressant d'explorer des méthodes d'amélioration continue du modèle, ainsi que l'intégration de fonctionnalités supplémentaires pour enrichir l'expérience utilisateur.

\par
